% Dot batch 81; source SHA-256: 6f102970d2dcc3026d1b373a2192d27429cdcb538b14203f7143a55b21ef59c9
% Archived source: dot_updates_20261124/batches/batch081/EGA第五卷_局部阶乘与线性系分类.tex
% Original body lines: 26--87; integration map: integration/late_integration.json
\begingroup
\def\newscope{V-L}
\def\depth{\operatorname{depth}}
\def\Spec{\operatorname{Spec}}
\def\Ass{\operatorname{Ass}}
\def\OO{\mathcal O}
\def\FF{\mathcal F}
\def\JJ{\mathcal J}
\egas{3.5}{推论}

沿用命题3.3的假设。设 \(X\) 的局部环都是阶乘环，例如 \(X\) 正则。
则 \(\mathcal D\to\mathcal U\) 是双射。特别地，若 \(X\) 是域 \(k\) 上局部有限型的正则预概形，
\(P\) 是 \(k\) 上的射影丛，则 \(X\) 上由 \(P^\vee\) 参数化的、无固定分支的除子线性系，与
满足命题3.1、3.3的非退化条件的有理映射 \(X\dashrightarrow P\) 一一对应。具体地，在几何纤维上，
逐个取关联点的闭包；其中每一个闭包的像都不得落在射影空间的某个超平面内。\cmt{L18-081-01}{法135手写把“\(X\) 正则”扩为局部环阶乘，法136手写补上不经超平面分解；英102将域上条件简写成整个映射不经超平面分解。按命题3.3对 \(\mathcal U\) 的定义，须对几何关联点逐一检验。整个像不落在某个超平面内，不能替代此条件，见本版核查报告 L18-081。}

事实上，\(X\) 的局部环若都是阶乘环，则 \(U\) 上任一可逆模都可延拓为 \(X\) 上的可逆模。
所以命题3.3 b) 所提条件自动满足。另一方面，由 Auslander--Buchsbaum 定理，正则局部环是阶乘环。\cmt{L18-081-02}{可核 \href{https://stacks.math.columbia.edu/tag/0AGB}{Stacks，0AGB}。原稿没有展开前一句的除子论证明；本版核查报告 L18-081 明确在局部诺特情形使用该方法，不把这一证明无条件推广到任意非诺特基底。法136大叉划删的后续段落，另见本版“第三篇删划草稿补译”中对应页段；其稿层身份及仍未辨的部分单列。}

\egasection{4\quad 除子线性系与可逆模}
使用第一小节的结果，我们现在要用 \(X\) 上的可逆层完整描述 \(X\) 上的线性系。
可以假定 \(Q=P^\vee\to S\) 是满射；记 \(\pi:X\times_S Q\to X\) 及 \(p:X\times_S Q\to Q\) 为两个投影。
按照原稿所引第20.3小节的一般结果，在 \(X\times_S Q\) 上给出一个有效除子 \(D\)，等价于给出一个可逆模 \(N\) 及它的一个正则截面 \(\phi\)（只计到保持截面的同构）。\cmt{L18-081-03}{法137保留第20.3小节的草稿引用；英104脚注78请作者确认是否为正式本21.3，未视为已经核定。有效 Cartier 除子与可逆层正则截面的对应可参见 \href{https://stacks.math.columbia.edu/tag/0C4S}{Stacks，0C4S}。英稿将“也满射”紧接该句；若指 \(p\)，仅由 \(Q\to S\) 满射不能推出，例如 \(X=\varnothing\)。若指 \(\pi\)，则由基变换成立。这里写明两投影且不要求 \(p\) 满射。}
说 \(D\) 是 \(X/S\) 上由 \(Q\) 参数化的除子线性系，就是说：

\noindent 1) 对每个 \(t\in Q\)，\(\phi\) 在 \(p\) 的纤维上诱导的截面 \(\phi_t\) 是正则截面；这也蕴含 \(\phi\) 正则。\cmt{L18-081-04}{这里“正则”即局部以非零因子表示，不要求处处非零。纤维正则给出相对有效 Cartier 除子，使用平坦且局部有限表示的环境及局部主闭子概形的纤维判据，见 \href{https://stacks.math.columbia.edu/tag/062Y}{Stacks，062Y(3)}。}

\noindent 2) 对每个 \(x\in X\)，\(N\) 在 \(\pi\) 的纤维上诱导的 \(N_x\) 具有度数一，即为 \(\OO(1)\) 型。\cmt{L18-081-05}{在正秩一层，射影纤维是点，\(\OO(1)\) 型仍有意义，但它不与不同次数的纤维线丛相区分。此处用上一小节和第一小节的一次约定，不能由点纤维的 Picard 群判别整数度数。}

由第1小节，在 \(X\) 上的射影丛 \(X\times_S Q\) 上给出满足 2) 的可逆模 \(N\)，等价于在 \(X\) 上给出可逆模 \(L\)。由 \(L\) 得到
\[
 N\simeq \pi^*L\otimes p^*\OO_Q(1),
\]
而由 \(N\) 恢复
\[
 L\simeq \pi_*(N(-1)),\qquad N(-1)=N\otimes p^*\OO_Q(-1).
\]
给出 \(\phi\)，就是给出 \(\pi^*L\otimes p^*\OO_Q(1)\) 的一个截面，也就是给出下述推前层的一个截面：
\[
 \pi_*\bigl(\pi^*L\otimes p^*\OO_Q(1)\bigr)
 \simeq L\otimes\pi_*p^*\OO_Q(1)
 \simeq L\otimes g^*(E^\vee).
\]
最后一个同构使用前面关于射影丛的计算；这里 \(g:X\to S\) 是结构态射，\(Q=\mathbf P_S(E^\vee)\)。因此，给出 \(\phi\) 等价于给出同态\cmt{L18-081-06}{法137的手写在右端补了 \(E_X^\vee\)，英104的转写漏上标。按一维商的约定，\(\pi_*p^*\OO_Q(1)=g^*(E^\vee)\)，不能写成 \(g^*E\)；后者是在把截面解释为同态时出现的。第一步用投影公式，在 \(L\) 局部平凡处即为恒等。射影丛约定见 \href{https://stacks.math.columbia.edu/tag/01OA}{Stacks，01OA}。}
\[
 g^*E\longrightarrow L,
 \qquad\text{或者等价地，}\quad u:E\longrightarrow g_*L.
\]


还需用 \(u\) 表述条件 1)。上述截面和同态 \(g^*E\to L\) 的构造与基变换相容，因而逐纤维考察即可。
若 \(K/k(s)\) 为扩域，则 \(Q_s\) 的 \(K\)-点对应于 \(E(s)\otimes_{k(s)} K\) 中的直线。
所以要求每个非零向量 \(t\in E(s)\otimes_{k(s)}K\) 所给出的截面 \(u(t)\) 在 \(X_K\) 上正则，其中 \(X_K=X_s\times_{k(s)}\Spec K\)。
只需在 \(k(s)\) 的一个代数闭包上检验。\cmt{L18-081-07}{法137末在射影语境下说“直线”，法138与英104--105略写为“每个 \(t\)”。必须排除零向量，且改取非零倍数不影响正则性。例如在非空整纤维上，零向量给出零截面而不正则。这里基变换相容的是截面及其伴随前的同态，不是断言任意 \(g_*L\) 的形成与任意基变换交换。}

\egas{4.1}{命题}

设 \(g:X\to S\) 平坦且局部有限表示，\(E\) 是 \(S\) 上处处具有正秩的有限型局部自由模，
\(P=\mathbf P_S(E)\)，\(Q=P^\vee=\mathbf P_S(E^\vee)\)。
则 \(X/S\) 上由 \(Q\) 参数化的除子线性系，与满足下述条件的序对 \((L,u)\) 的同构类一一对应：
\(L\) 是 \(X\) 上的可逆模，\(u:E\to g_*L\) 是同态。对每个 \(s\in S\)，每个扩域 \(K/k(s)\)，
以及每个非零 \(t\in E(s)\otimes_{k(s)}K\)，对应的截面 \(u(t)\) 在 \(X_K\) 上正则。
可只取 \(K=\overline{k(s)}\)。\cmt{L18-081-08}{命题中法、英均出现“\(X\) 上”的射影丛，按前面满射 \(P^\vee\to S\) 及 \(u:E\to g_*L\) 的型别订为“\(S\) 上”。法138的边注补正秩，英稿也明确非空纤维。序对的同构是 \(L\simeq L'\) 且借此带 \(u\) 为 \(u'\)，不是对 \(E\) 另换坐标。}

若 \(X\to S\) 的纤维几何整，则 \(u\) 的条件简化为：对所有 \(s\in S\)，自然映射
\[
 u(s):E(s)\longrightarrow H^0(X_s,L_s)
\]
是单射。作者建议在4.1中明说此点，并为便于查阅补记 \(D\) 的构造：由 \(u\) 得到
\(\pi^*L\otimes p^*\OO_Q(1)\) 的截面 \(\phi\)，取其零除子。
这里 \(u(s)\) 由 \(g^*E\to L\) 限制到纤维后再取伴随。\cmt{L18-081-09}{几何整保证几何纤维上的非零截面正则。对域扩张，纤维是拟紧拟分离时，全局截面直接与平坦扩域交换；对仅局部有限表示的一般纤维，可用有限维 \(E(s)\) 在有限多个仿射开上检验单射，见本版核查报告 L18-081。不能径用任意 \(g_*L\) 的任意基变换同构。}
\par\endgroup
